\cleardoublepage
\chapter{ANALISIS KEBUTUHAN}
\label{chap:analisis-kebutuhan}

Bab ini menguraikan proses penggalian kebutuhan pengembangan ADADeH. Pembahasan diawali dengan metode pengumpulan data pada Subbab \ref{metode_data}, dilanjutkan dengan analisis kebutuhan pengguna pada Subbab \ref{keb_pengguna}. Hasil kedua subbab tersebut menjadi dasar pemilihan alternatif solusi dan perumusan kebutuhan sistem pada subbab berikutnya.

\section{Metode Pengumpulan Data}
\label{metode_data}
{Belum BERES}

\section{Analisis Kebutuhan Pengguna}
\label{keb_pengguna}

Berdasarkan data yang terkumpul pada Subbab \ref{metode_data}, kebutuhan pengguna ADADeH dirumuskan pada subbab ini. Kebutuhan pengguna adalah kemampuan yang harus dimiliki sistem dari sudut pandang pengguna, belum menyangkut cara pemenuhannya secara teknis.

\subsection{Identifikasi Pengguna}
\label{identifikasi_pengguna}

ADADeH ditujukan bagi dua kelompok pengguna berikut.
\begin{enumerate}[(a).]
  \item \textbf{Pengguna awam.} Kelompok ini merupakan pengguna utama, yaitu korban atau calon korban kejahatan berbasis \textit{deepfake audio}. Kelompok ini tidak memiliki latar belakang pemrosesan sinyal maupun forensik digital, sehingga membutuhkan keluaran yang langsung dapat dipahami.
  \item \textbf{Pendamping pelaporan.} Kelompok ini mencakup pihak yang membantu korban menyiapkan laporan. Kelompok ini membutuhkan keluaran yang dapat disimpan atau dilampirkan sebagai indikasi awal, dengan kesadaran bahwa hasil tersebut bukan pengganti pemeriksaan ahli.
\end{enumerate}

\subsection{Rumusan Kebutuhan Pengguna}
\label{rumusan_kp}

Kebutuhan pengguna diberi kode KP-01 sampai KP-08 dan diuraikan sebagai berikut.

\begin{enumerate}[{KP-}01.]
  \item \textbf{Pengguna membutuhkan sarana pemeriksaan keaslian rekaman yang dapat diakses secara mandiri.}\\
  Survei menunjukkan sebagian besar responden tidak mengetahui prosedur pembuktian, sementara laboratorium forensik digital jumlahnya terbatas \autocite{SIP_Law_Firm_2025}. Pengguna karena itu membutuhkan sistem yang dapat dipakai tanpa melalui prosedur kelembagaan terlebih dahulu.

  \item \textbf{Pengguna membutuhkan sistem yang menerima rekaman dari perangkat sehari-hari.}\\
  Rekaman pengguna umumnya berasal dari ponsel dan tersimpan dalam format \textit{lossy}. Sistem harus menerima berkas tersebut dan menanganinya sendiri, tanpa menuntut pengguna melakukan konversi format secara manual.

  \item \textbf{Pengguna membutuhkan keluaran berupa kesimpulan yang langsung dapat dipahami.}\\
  Responden mengharapkan pernyataan asli atau \textit{deepfake} beserta tingkat keyakinannya, bukan keluaran teknis berupa nilai \textit{logit} atau grafik spektrum. Kebutuhan ini sejalan dengan penyajian skor keyakinan $c = \max(p,\, 1-p)$ pada Subbab \ref{sist_mulmed}.

  \item \textbf{Pengguna membutuhkan keterangan mengenai keterbatasan hasil.}\\
  Karena keluaran menyangkut kepentingan hukum \autocite{Iman_2026}, sistem harus menegaskan bahwa hasil merupakan indikasi awal dan tidak menggantikan pemeriksaan forensik digital oleh ahli.

  \item \textbf{Pengguna membutuhkan hasil dalam waktu yang singkat.}\\
  Pemeriksaan umumnya dilakukan pada situasi mendesak, misalnya saat pengguna sedang dihubungi pelaku. Waktu tunggu yang lama membuat sistem kehilangan manfaatnya.

  \item \textbf{Pengguna membutuhkan sistem yang dapat memeriksa rekaman berdurasi panjang.}\\
  Rekaman percakapan dapat berlangsung beberapa menit, sedangkan model bekerja pada potongan berdurasi tetap. Pengguna membutuhkan satu kesimpulan menyeluruh atas rekaman yang diunggah, bukan daftar hasil per potongan.

  \item \textbf{Pengguna membutuhkan peringatan ketika mutu rekaman tidak memadai.}\\
  Rekaman yang terlalu singkat atau mengandung derau tinggi menurunkan ketepatan deteksi. Pengguna perlu mengetahui hal tersebut agar tidak menyimpulkan secara keliru dari hasil yang kurang dapat dipercaya.

  \item \textbf{Pengguna membutuhkan jaminan kerahasiaan atas rekaman yang diunggah.}\\
  Rekaman suara merupakan data pribadi dan kerap memuat percakapan sensitif. Pengguna membutuhkan kepastian bahwa berkas tidak disebarluaskan dan tidak disimpan melebihi keperluan pemrosesan.
\end{enumerate}

\subsection{Keterlacakan Kebutuhan}
\label{keterlacakan}

Agar setiap kebutuhan dapat ditelusuri asalnya, Tabel \ref{tab:keterlacakan} memetakan kebutuhan pengguna terhadap sumber data dan rumusan masalah pada Bab \ref{chap:pendahuluan}. Pemetaan ini sekaligus memastikan tidak ada kebutuhan yang dirumuskan tanpa dasar.

\begin{table}[t]
  \centering
  \caption{Keterlacakan kebutuhan pengguna terhadap sumber data dan rumusan masalah}
  \label{tab:keterlacakan}
  \begin{tabularx}{\textwidth}{| c | X | c |}
    \hline
    \textbf{Kode} & \textbf{Sumber Data} & \textbf{Rumusan Masalah} \\
    \hline
    KP-01 & Studi literatur dan survei butir 5 & 3 \\
    KP-02 & Karakterisasi data dan survei butir 8 & 2 \\
    KP-03 & Survei butir 3 dan butir 7 & 3 \\
    KP-04 & Studi literatur & 3 \\
    KP-05 & Survei butir 4 & 2 \\
    KP-06 & Karakterisasi data dan survei butir 6 & 2 \\
    KP-07 & Karakterisasi data dan studi literatur & 1 \\
    KP-08 & Studi literatur & 3 \\
    \hline
  \end{tabularx}
\end{table}

Kedelapan kebutuhan tersebut menjadi dasar penyusunan alternatif solusi pada subbab berikutnya. Setiap alternatif dinilai berdasarkan kemampuannya memenuhi KP-01 sampai KP-08, kemudian alternatif terpilih diterjemahkan menjadi kebutuhan sistem yang bersifat teknis.

\section{Pemilihan Alternatif Solusi}
\label{metode_data}
{Siapa yg mau ambil}

\section{Analisis Kebutuhan Sistem}
\label{metode_data}
{Siapa yg mau ambil}